% Generated draft: requires mathematical and source review.
\documentclass[11pt]{article}
\usepackage[margin=25mm]{geometry}
\usepackage{fontspec}
\setmainfont{Latin Modern Roman}
\usepackage{amsmath,amssymb,mathtools,unicode-math}
\setmathfont{Latin Modern Math}
\usepackage{enumitem,needspace,xurl,hyperref}
\setlength{\parindent}{0pt}
\setlength{\parskip}{5pt}
\setlength{\emergencystretch}{3em}
\newcommand{\N}{\mathbb{N}}
\newcommand{\Z}{\mathbb{Z}}
\newcommand{\Q}{\mathbb{Q}}
\newcommand{\R}{\mathbb{R}}
\newcommand{\C}{\mathbb{C}}
\newcommand{\F}{\mathbb{F}}
\newcommand{\A}{\mathbb{A}}
\newcommand{\PP}{\mathbb{P}}
\newcommand{\E}{\mathbb{E}}
\newcommand{\eps}{\varepsilon}
\newcommand{\dd}{\,\mathrm{d}}
\begin{document}
\section*{UnsolvedMath ID: 1100218}
\textbf{Source problem code:} catalogue entry
\subsection{Definitions and mathematical statement}
A metric space $M$ is proper if its closed bounded balls are compact. It is CAT$(0)$ if it is geodesic and distances in each geodesic triangle are bounded above by the corresponding distances in its Euclidean comparison triangle. The action of $\operatorname{Isom}(M)$ is cocompact if there is a compact set $K\subseteq M$ with $\operatorname{Isom}(M)K=M$.

An infinite geodesic ray from $a$ is an isometric embedding $\gamma:[0,\infty)\to M$ with $\gamma(0)=a$. The source defines $M$ to be \emph{almost geodesically complete} if
\[
\exists R\geq0\ \forall a,b\in M\ \exists\gamma:[0,\infty)\to M\ \exists t\geq0:
\quad \gamma(0)=a,\quad
 d(\gamma(s),\gamma(u))=|s-u|\ (s,u\geq0),\quad
 d(\gamma(t),b)\leq R.
\]
The constant $R$ is uniform in $a,b$, while the ray may depend on both points.

Ross Geoghegan's question as transcribed asks whether every proper CAT$(0)$ space $M$ with cocompact isometry group has this property. The transcribed question omits a noncompactness hypothesis. The supplied reported affirmative theorem instead states:
\[
\text{If $M$ is noncompact, proper, CAT$(0)$, and $\operatorname{Isom}(M)$ acts cocompactly, then $M$ is almost geodesically complete.}
\]
These scopes must be distinguished: a nonempty compact metric space cannot contain an isometrically embedded infinite ray, so it cannot satisfy the displayed definition. Noncompactness cannot silently be removed from the cited affirmative statement or inserted into the historical wording.

The supplied research summary attributes the noncompact assertion to Theorem 5 of Geoghegan--Ontaneda (2007), combining nonvanishing compactly supported integral cohomology $H_c^d(M;\Z)$ in a top degree $d$ with Ontaneda's earlier homological criterion. It also reports that the special case of CAT$(0)$ groups, acting geometrically with discrete orbits, was already settled by Ontaneda's Theorem B (2005). These literature statements are supplied assessments, rather than independently verified proofs here.
\subsection{Short English statement}
Almost geodesic completeness means that every ordered pair of points can be connected up to one uniform error by an infinite geodesic ray starting at the first point. The supplied literature reports this for noncompact proper CAT$(0)$ spaces with cocompact isometry group. The source transcription lacks noncompactness, which is incompatible with compact spaces and requires correction or qualification.
\subsection{Sources}
\begin{itemize}
\item Local UnsolvedMath registry, permanent ID 1100218, exact source code catalogue entry, ``Questions in Geometric Group Theory --- Q 2.18''; registry status solved and classification SOLVED-IN-LITERATURE.
\item Bestvina, \emph{Questions in Geometric Group Theory} (2004), Question 2.18, PDF page 9, attributed to Ross Geoghegan, \url{https://www.math.utah.edu/~bestvina/eprints/questions-updated.pdf}. The supplied record identifies manual transcription of a public author-hosted PDF, accessed July 29, 2026, and says the author PDF records partial progress or special cases, with current status still requiring release review.
\item Geoghegan--Ontaneda, \emph{Topology} 46 (2007), 129--137, Theorem 5, cited by the supplied research summary for the affirmative noncompact proper cocompact CAT$(0)$ result and its compactly supported cohomology input. The record provides no paper title.
\item Ontaneda, \emph{Topology} 44 (2005), 47--62, cited for the homological criterion and, specifically, Theorem B for the CAT$(0)$-group case. The record provides no paper title. The research summary expressly makes no claim of an independent solution.
\end{itemize}
\end{document}
